\section{Conclusões}
\label{conclusoes}

O trabalho partiu de uma constatação do levantamento bibliográfico: apesar de amplamente usado, o FedProx não tem uma vantagem consistente sobre o FedAvg na literatura recente, e os estudos raramente controlam $\mu$, o nível de heterogeneidade ou o custo no cliente. Para responder em quais condições o termo proximal compensa, foi desenhado um experimento pareado, reprodutível e de custo zero, que cruza tipo e nível de heterogeneidade com o valor de $\mu$ e mede tanto a qualidade do modelo quanto o custo local. \pendente{Concluir com a resposta às QP1 e QP2 em uma ou duas frases, com base nos resultados.}
